\documentclass[11pt]{article}
\usepackage[margin=1in]{geometry}
\usepackage{amsmath,amssymb}
\usepackage{array}
\usepackage{booktabs}
\usepackage{longtable}
\usepackage{tabularx}
\usepackage{xcolor}
\usepackage[hidelinks]{hyperref}

\newcolumntype{L}[1]{>{\raggedright\arraybackslash}p{#1}}
\newcolumntype{Y}{>{\raggedright\arraybackslash}X}
\setlength{\LTpre}{0.4em}
\setlength{\LTpost}{0.8em}
\setlength{\emergencystretch}{3em}
\hbadness=10000
\hfuzz=30pt

\newcommand{\code}[1]{\texttt{\detokenize{#1}}}
\newcommand{\vcmd}{v^{\mathrm{cmd}}_t}
\newcommand{\vcorr}{v^{\mathrm{corr}}_t}
\newcommand{\vgov}{\bar v^g_t}
\newcommand{\vfin}{\tilde v_t}
\newcommand{\vbrake}{v^{\mathrm{brake}}_t}
\newcommand{\vsafe}{v_{\mathrm{safe}}}
\newcommand{\amax}{a_{\max}}

\title{\textbf{instinctRL Development Handbook v1.1:\\Platform-Locked Code-Aware Implementation Plan}}
\author{NavRL / instinctRL}
\date{}

\begin{document}
\maketitle
\tableofcontents
\newpage

%============================================================
\section{Title and Purpose}

This handbook is the platform-locked code-aware development plan for
\textbf{instinctRL}. The scientific target remains Paper~1, derived from
\code{docs/paper1\_vel\_ctrl.tex} and the method-aligned v0 handbook
\code{docs/paper1\_vel\_ctrl\_handbook.tex}. This document does not redesign
the method or change the scientific positioning.

The locked method is:
\begin{itemize}
\item range--inertial measurement-space velocity governor;
\item learned governor outputs a corrected body-frame velocity command;
\item velocity-controller execution, not CTBR, motor thrust, body-rate, or
attitude-level learned action;
\item measurement-space anchor for station-keeping;
\item range-Jacobian observability logging for analysis/evaluation;
\item ICS-inspired command attenuation, not a formal CBF guarantee.
\end{itemize}

\paragraph{Development name.}
The code and implementation project name is \textbf{instinctRL}. Module names,
config keys, stages, tickets, logs, and implementation artifacts must use
\code{instinctRL}. The term Paper~1 may be used only to refer to the scientific
target and paper route.

\paragraph{Platform and sensor lock.}
The LiDAR sensor is \textbf{Livox MID360}. The UAV platform is
\textbf{TASLAB\_UAV}. The simulator already contains MID360 simulation helpers
and a TASLAB\_UAV OmniDrones model/controller interface; instinctRL must reuse
these assets. A generic LiDAR, generic ray sensor, generic drone, or new
simulator may be used only as a debugging-only fallback and must be labeled as
such in configs and logs.

\paragraph{Hard information contract.}
No pose, odometry, explicit velocity estimate, reconstructed map, SLAM state,
or privileged simulator state may enter the deployed actor input. Privileged
state may be used only for reward, critic input, evaluation, logging, or
upper-bound baselines.

\paragraph{Domain note.}
The repository root \code{CONTEXT.md} still describes a broader INSTINCT route
with CTBR, required recurrent actors, CMDPs, and PPO-Lagrangian. That glossary
does not override this platform-locked instinctRL handbook. For this document,
\code{paper1\_vel\_ctrl.tex} is authoritative.

%============================================================
\section{Platform and Sensor Infrastructure Audit}

\subsection{instinctRL-0: Platform and Sensor Infrastructure Audit}

This stage must be completed before new learning modules are implemented. Its
purpose is to prove that the locked MID360--TASLAB\_UAV platform stack is
available, correctly wired, and actor-input compliant.

\begin{longtable}{@{}L{0.22\textwidth}L{0.36\textwidth}L{0.32\textwidth}@{}}
\caption{instinctRL-0 acceptance criteria.}\\
\toprule
\textbf{Audit item} & \textbf{Repository evidence to inspect/reuse} & \textbf{Acceptance criterion} \\
\midrule
\endfirsthead
\toprule
\textbf{Audit item} & \textbf{Repository evidence to inspect/reuse} & \textbf{Acceptance criterion} \\
\midrule
\endhead
MID360 ray pattern and frame convention & \code{training/envs/livox\_mid360.py}: \code{LivoxMid360Config}, \code{LivoxMid360Pattern}, \code{generate\_rays()}, \code{create\_livox\_from\_hydra\_cfg()} & Stable ray ordering, documented body-frame convention \(X\) forward, \(Y\) left, \(Z\) up, and deterministic RL mode available. \\
MID360 range output & \code{NavigationEnv} RayCaster path in \code{training/scripts/env.py}; integration helper in \code{training/scripts/livox\_mid360\_integration.py}; tests in \code{training/unit\_test/test\_livox\_mid360.py} & MID360 produces a flat range vector \(r_t\) with timestamps and consistent ray count. \\
Mask and dropout handling & \code{LivoxMid360Pattern.apply\_noise()} already models Gaussian noise, dropout, near-range dropout, and invalid \(\infty\) returns. & Valid-return mask \(m_t\) is derived from finite in-range MID360 returns; dropout is represented explicitly. \\
Reliability weights & Neighbor consistency can reuse \code{training/envs/lidar\_processor.py} projection/debug metadata and MID360 validity/noise signals. & Reliability weights \(w_t\in[0,1]\) are computed without surface normals and logged. \\
Update rate and timestamps & MID360 file documents 10 Hz frame rate; active simulation calls \code{lidar.update(self.dt)} in \code{NavigationEnv.\_post\_sim\_step}. & Logs include sensor timestamp/frame age; history buffer marks stale or repeated frames. \\
TASLAB\_UAV model & \code{omni\_drones/robots/drone/taslab\_uav.py}, \code{assets/usd/taslab\_uav.yaml}, \code{taslab\_uav.usd}, \code{drone/\_\_init\_\_.py}. & \code{MultirotorBase.REGISTRY["TaslabUAV"]} resolves and spawns the locked platform. \\
TASLAB\_UAV parameters & \code{taslab\_uav.yaml}; \code{docs/TASLAB\_UAV\_README.md}; \code{unit\_test/test\_hover.py}, \code{test\_flight.py}. & Mass, inertia, rotor layout, force constants, and max rotor speeds are loaded from TASLAB\_UAV parameters. \\
Velocity controller interface & \code{train.py} uses \code{VelController(LeePositionController)}; TASLAB gains in \code{lee\_controller\_taslab\_uav.yaml}. & TASLAB\_UAV accepts corrected body-frame velocity commands through an adapter into the existing velocity-controller path. \\
Actor-input contract & Current \code{env.py}, \code{ppo.py}, and ROS wrappers. & No pose, odometry, explicit velocity, map, or privileged simulator state reaches deployed actor inputs. \\
\bottomrule
\end{longtable}

\paragraph{Known platform-lock gaps in current code.}
The active \code{NavigationEnv} creates a RayCaster with
\code{prim\_path="/World/envs/env_.*/Hummingbird\_0/base\_link"} even though
\code{training/cfg/drone.yaml} selects \code{drone.model\_name: "TaslabUAV"}.
This must be fixed by resolving the spawned TASLAB\_UAV base link dynamically
or by formatting the prim path from \code{cfg.drone.model\_name}. The active
environment also uses \code{patterns.BpearlPatternCfg} directly; MID360 helpers
exist, but instinctRL-0 must verify whether the active training path is using
the MID360 configuration, ray ordering, masks, dropout, and timestamps.

%============================================================
\section{Repository Overview}

\begin{longtable}{@{}L{0.16\textwidth}L{0.25\textwidth}L{0.20\textwidth}L{0.18\textwidth}L{0.15\textwidth}@{}}
\caption{Existing infrastructure and instinctRL reuse plan.}\\
\toprule
\textbf{Subsystem} & \textbf{Existing files / directories} & \textbf{Current role} & \textbf{instinctRL reuse} & \textbf{Missing pieces} \\
\midrule
\endfirsthead
\toprule
\textbf{Subsystem} & \textbf{Existing files / directories} & \textbf{Current role} & \textbf{instinctRL reuse} & \textbf{Missing pieces} \\
\midrule
\endhead
Training entry points & \code{isaac-training/training/scripts/train.py}, \code{eval.py} & Launch Isaac Sim, build \code{NavigationEnv}, wrap with \code{VelController}, train/evaluate PPO, log WandB. & Reuse entry points; add \code{instinctRL.enabled} mode and platform audit. & instinctRL metrics, input audit, platform-lock checks. \\
Environment & \code{training/scripts/env.py} & Builds TASLAB-configured drone, terrain, current RayCaster, observations, rewards, done flags. & Reuse \code{NavigationEnv}; gate instinctRL observation/reward/action path. & Hardcoded Hummingbird LiDAR path; actor leaks pose/velocity. \\
MID360 simulation & \code{training/envs/livox\_mid360.py}; \code{training/scripts/livox\_mid360\_integration.py}; \code{unit\_test/test\_livox\_mid360.py}. & MID360 config, ray pattern, mount transform, occlusion, noise/dropout, helper tests. & Reuse; do not create a new sensor simulator. & Wire helper into active \code{NavigationEnv}; add timestamps/reliability output. \\
Range projection & \code{training/envs/lidar\_processor.py}. & Converts point clouds to body-frame depth images with grid masks/debug info. & Reuse for optional MID360 range image/CNN representation and neighbor consistency. & Preserve raw range vector for anchor/ICS. \\
TASLAB\_UAV model & \code{third\_party/OmniDrones/.../taslab\_uav.py}, \code{taslab\_uav.yaml}, \code{taslab\_uav.usd}. & Registered OmniDrones multirotor with TASLAB parameters. & Reuse as locked platform. & Verify prim path and controller gains in active training. \\
Velocity controller & \code{VelController}, \code{LeePositionController}; \code{lee\_controller\_taslab\_uav.yaml}. & Converts velocity command to low-level rotor action. & Reuse; instinctRL stays velocity-controller-based. & Body-frame to controller-frame adapter and audit logs. \\
Observation pipeline & \code{env.py:\_set\_specs}, \code{\_compute\_state\_and\_obs}. & Current actor sees \code{lidar}, \code{state}, \code{direction}, \code{dynamic\_obstacle}. & Reuse only the range acquisition and privileged logging parts. & Remove actor pose/velocity/map leaks; add history. \\
Policy network & \code{training/scripts/ppo.py}. & CNN lidar encoder, dynamic-obstacle MLP, Beta actor outputting raw velocity. & Reuse PPO loop and parts of encoder. & Governor head \((\alpha_t,\vcorr)\), actor/critic split. \\
Reward code & Inline in \code{env.py}. & Goal-tracking velocity reward, safety, smoothness, height, collision. & Reuse simulator quantities reward-only. & instinctRL tracking, anchor, intervention, safety terms. \\
Command generation & \code{training/scripts/command\_generator.py}. & Vectorized adversarial command modes. & Reuse for training/evaluation command regimes only. & Feed only \(\vcmd\), not privileged generator inputs, to actor. \\
Logging/evaluation & \code{train.py}, \code{utils.py}, \code{EpisodeStats}, WandB. & Basic return/reach/collision/video logging. & Reuse logging plumbing. & Anchor, ICS, observability, audit, platform-lock logs. \\
ROS deployment & \code{ros1/navigation\_runner/scripts/navigation.py}, \code{policy\_server.py}, \code{srv/*.srv}, \code{safeAction.cpp}. & Current deployment uses odometry, map raycast, dynamic-obstacle services, ORCA-style safety. & Reuse only process/service/publisher scaffolding. & New MID360-range and IMU policy request excluding odom/map actor inputs. \\
\bottomrule
\end{longtable}

%============================================================
\section{instinctRL Module-to-Code Mapping}

\begin{longtable}{@{}L{0.08\textwidth}L{0.16\textwidth}L{0.21\textwidth}L{0.25\textwidth}L{0.20\textwidth}@{}}
\caption{instinctRL modules, reuse first.}\\
\toprule
\textbf{ID} & \textbf{Method role} & \textbf{Existing code to reuse} & \textbf{Changes or new functions} & \textbf{Why new code is needed} \\
\midrule
\endfirsthead
\toprule
\textbf{ID} & \textbf{Method role} & \textbf{Existing code to reuse} & \textbf{Changes or new functions} & \textbf{Why new code is needed} \\
\midrule
\endhead
instinctRL-M0 & Platform and sensor audit. & MID360 files, TASLAB files, \code{train.py}, \code{env.py}, unit tests. & Add \code{instinctRL/platform\_audit.py} or test-only audit helper. & No current single check proves MID360 is attached to TASLAB and actor-clean. \\
instinctRL-M1 & Observation/history buffer. & \code{env.py} range extraction; \code{livox\_mid360.py}; \code{lidar\_processor.py}. & Add history buffer and actor schema under \code{instinctRL/observation.py}. & Existing observation contains forbidden pose/velocity and has no history/timestamps. \\
instinctRL-M2 & Measurement-space anchor. & Reuse MID360 \(r,m,w\) from M1. & Add \code{instinctRL/anchor.py}: hysteresis, capture, masked error, logs. & No existing anchor lifecycle or anchor loss exists. \\
instinctRL-M3 & Velocity governor network. & Reuse PPO loop and LiDAR encoder from \code{ppo.py}. & Add governor head or \code{instinctRL/governor.py} producing \(\alpha_t,\vcorr\). & Current actor outputs raw velocity and consumes forbidden fields. \\
instinctRL-M4 & ICS-inspired attenuation. & Reuse range/history from M1 and command path from \code{train.py}. & Add \code{instinctRL/ics.py}: range-rate fit, active set, \(\beta_t\). & Existing \code{safeAction.cpp} is ORCA/map/obstacle-state based and not actor-contract compliant. \\
instinctRL-M5 & Observability logger. & Reuse RayCaster scene and evaluation logger. & Add \code{instinctRL/observability.py}. & No existing range-Jacobian or drift correlation logger. \\
instinctRL-M6 & Reward terms. & Reuse privileged state and collision computation in \code{env.py}. & Add \code{instinctRL/rewards.py} or gated reward functions. & Existing reward is goal navigation, not measurement-space station-keeping. \\
instinctRL-M7 & Evaluation/logging/input audit. & Reuse \code{utils.evaluate}, \code{EpisodeStats}, WandB. & Add \code{instinctRL/audit.py}, metrics hooks. & Existing logging lacks actor contract and Paper~1 contribution metrics. \\
instinctRL-M8 & Baseline/ablation harness. & Reuse configs, command generator, train/eval scripts. & Add config switches \code{instinctRL.baseline.id}. & Existing configs do not isolate B0--B8. \\
instinctRL-M9 & Real-robot deployment wrapper. & Reuse ROS nodes/services as scaffolding. & Add MID360/IMU policy wrapper and message schema. & Current ROS actor input path uses odometry/map raycast/obstacle state. \\
\bottomrule
\end{longtable}

%============================================================
\section{Actor Input Contract Audit}

\begin{longtable}{@{}L{0.23\textwidth}L{0.22\textwidth}L{0.39\textwidth}@{}}
\caption{Current and target actor inputs.}\\
\toprule
\textbf{Candidate field} & \textbf{Current source} & \textbf{instinctRL status} \\
\midrule
\endfirsthead
\toprule
\textbf{Candidate field} & \textbf{Current source} & \textbf{instinctRL status} \\
\midrule
\endhead
MID360 range \(r_t\) & RayCaster hit distance or \code{LivoxMid360Pattern} generated rays. & Allowed actor input. Must be true distance, not only danger-coded inverse range. \\
Valid-return mask \(m_t\) & Derived from finite in-range MID360 returns and dropout. & Allowed actor input. Required. \\
Reliability weights \(w_t\) & Derived from validity, dropout, staleness, neighbor consistency. & Allowed actor input. Required; no surface normals. \\
Ray directions / ordering metadata & MID360 pattern in body frame. & Allowed for deterministic ICS and preprocessing; actor may receive only configured representation. \\
IMU / attitude cues & TASLAB/drone state or ROS IMU. & Allowed only as derived body angular rate, gravity direction, and attitude cue; no translational pose/velocity. \\
Operator command \(\vcmd\) & Operator or command generator output. & Allowed actor input in body/governor frame. \\
Previous issued command & New governor state. & Allowed actor input. \\
History \(h_t\) & New buffer over allowed fields. & Allowed actor input. \\
\code{state}: \code{rpos\_clipped\_g}, \code{distance\_2d}, \code{distance\_z} & \code{env.py} from target and drone position. & Forbidden deployed actor input. Keep only for reward/evaluation if needed. \\
\code{vel\_g} & \code{env.py} from root-state velocity. & Forbidden deployed actor input. Reward/critic/evaluation only. \\
\code{direction} & \code{env.py}/\code{ppo.py} target direction. & Forbidden for deployed actor. Replace goal-frame policy with body-frame command convention. \\
\code{dynamic\_obstacle} & Simulator obstacle pose/velocity relative to drone. & Forbidden deployed actor input. Evaluation/upper-bound only. \\
\code{info/drone\_state} and root state & \code{env.py}. & Privileged; reward/critic/logging only. \\
ROS odometry/map raycast & \code{navigation.py}. & Forbidden deployed actor input. May be used for low-level infrastructure or evaluation only. \\
\bottomrule
\end{longtable}

\paragraph{Audit implementation.}
Add an \code{instinctRL.audit} check that runs at environment construction,
policy initialization, rollout collection, evaluation, checkpoint export, and
ROS inference. It must fail on keys or producer metadata containing
\code{pose}, \code{pos}, \code{position}, \code{odom}, \code{velocity},
\code{vel\_g}, \code{vel\_w}, \code{root\_state}, \code{map}, \code{slam},
\code{privileged}, \code{direction}, \code{distance}, or
\code{dynamic\_obstacle}, unless the field is explicitly in a reward-only,
critic-only, logging-only, or baseline-upper-bound branch.

%============================================================
\section{Action Interface Plan}

The governor action remains velocity-controller based:
\[
  \vgov = \alpha_t\vcmd + \vcorr,\qquad
  \vfin = \beta_t\vgov + (1-\beta_t)\vbrake .
\]

\begin{longtable}{@{}L{0.28\textwidth}L{0.62\textwidth}@{}}
\caption{Velocity-controller action path on TASLAB\_UAV.}\\
\toprule
\textbf{Step} & \textbf{Implementation plan} \\
\midrule
\endfirsthead
\toprule
\textbf{Step} & \textbf{Implementation plan} \\
\midrule
\endhead
Governor output & Reuse \code{ppo.py} training loop but replace actor head so it outputs \(\alpha_t\in[0,1]\) and bounded \(\vcorr\). \\
Body-frame convention & Define \(\vcmd,\vcorr,\vgov,\vfin\) in TASLAB body/governor frame. Use attitude/yaw information only in the command adapter, not as translational actor state. \\
Clipping and smoothing & Clip \(\vcorr\), clip \(\vgov\) to velocity bounds, and smooth/rate-limit before ICS. Store previous issued command in the history buffer. \\
ICS attenuation & Apply \(\beta_t\) in \code{instinctRL/ics.py} after smoothing and before the controller transform. Empty active set returns \(\beta_t=1\). \\
Velocity controller execution & Reuse \code{VelController(LeePositionController)} and TASLAB parameters. The adapter converts the final body-frame velocity command to the velocity frame expected by \code{VelController}. \\
Forbidden interfaces & Do not expose TASLAB rotor actions, CTBR, body rates, motor thrust, or attitude setpoints as the learned instinctRL action. \\
\bottomrule
\end{longtable}

%============================================================
\section{Observation and History Buffer Implementation Plan}

\begin{longtable}{@{}L{0.25\textwidth}L{0.61\textwidth}@{}}
\caption{Observation mapping for MID360--TASLAB\_UAV.}\\
\toprule
\textbf{Term} & \textbf{Concrete plan} \\
\midrule
\endfirsthead
\toprule
\textbf{Term} & \textbf{Concrete plan} \\
\midrule
\endhead
\(r_t\) & Use MID360 ray ordering from \code{LivoxMid360Pattern} or the active RayCaster configured with MID360 angles. Compute \(r_i=\|p^{hit}_i-p^{lidar}\|\), not only \code{lidar\_range-r_i}. \\
\(m_t\) & Use finite, in-range, non-occluded, non-dropout returns. Invalid and dropout beams become mask zero. \\
\(w_t\) & Combine valid return, staleness, noise/dropout confidence, and neighbor consistency. Reuse \code{lidar\_processor.py} projection where helpful. \\
Timestamps & Record MID360 frame time, simulation time, and age relative to policy step. Mark stale or repeated frames. \\
IMU/attitude cues & Derive body angular velocity and gravity direction from TASLAB state or ROS IMU; keep position and translational velocity out of actor input. \\
\(\vcmd\) & Source from operator or \code{AdversarialCommandGenerator}; generator may use privileged information internally for scenario generation, but only the command enters the actor. \\
Previous action & Store \(\tilde v_{t-1}\) and optionally \(\bar v^g_{t-1}\), bounded and frame-labeled. \\
History \(h_t\) & Maintain a fixed window over allowed fields, masks, weights, timestamps, and previous command. \\
\bottomrule
\end{longtable}

Tests: MID360 shape/ray count, ray-order stability, timestamp monotonicity,
mask/dropout behavior, reliability bounds, history rollover, stale-frame
handling, and actor-input absence of pose/odometry/velocity/map/privileged
state.

%============================================================
\section{Measurement-Space Anchor Implementation Plan}

The anchor manager is new because no current code captures or regulates a
measurement-space station-keeping reference. It must consume only \(r_t,m_t,w_t\)
and \(\vcmd\).

\begin{longtable}{@{}L{0.28\textwidth}L{0.58\textwidth}@{}}
\caption{Anchor requirements.}\\
\toprule
\textbf{Function} & \textbf{Implementation} \\
\midrule
\endfirsthead
\toprule
\textbf{Function} & \textbf{Implementation} \\
\midrule
\endhead
Null-command hysteresis & \(\chi_t\) activates when \(\|\vcmd\|\le\epsilon_0\), deactivates when \(\|\vcmd\|>\epsilon_1\), with \(\epsilon_1>\epsilon_0\). \\
Anchor capture & On activation rising edge, freeze MID360 \(r^\star,m^\star\) and timestamp. \\
Anchor reset & Reset on episode reset, explicit operator reset, large-command exit, or insufficient valid anchor fraction. \\
Masked anchor error & \(e^r_t=(m_t\odot m^\star\odot w_t)\odot(r_t-r^\star)\). \\
Anchor loss & Robust masked loss active only while \(\chi_t=1\). No pose or odometry. \\
Logs & Active fraction, activation count, hold duration, valid-anchor fraction, mean/max anchor error, reset reason. \\
\bottomrule
\end{longtable}

%============================================================
\section{Governor Network Implementation Plan}

Reuse the PPO collection/update loop in \code{training/scripts/ppo.py}. Reuse
the existing LiDAR CNN only if the MID360 representation is image-like; if the
canonical actor input is a flat MID360 vector with masks/weights/history, add a
small history-stacking MLP or TCN encoder. GRU/LSTM may be added only as an
ablation.

\begin{itemize}
\item Actor input: allowed MID360 range/mask/weights, IMU cues, command,
previous command/output, anchor features, and history.
\item Actor output: \(\alpha_t\) and \(\vcorr\), followed by deterministic
formation of \(\vgov\).
\item Critic input: may include privileged TASLAB pose/velocity/clearance only
through a separate critic-only branch.
\item Configs: \code{instinctRL.governor.encoder}, \code{alpha\_mode},
\code{v\_corr\_limit}, \code{velocity\_limit}, \code{smoothing\_tau}.
\item Tests: output bounds, no critic-to-actor leakage, export compatibility,
and deterministic mean action.
\end{itemize}

%============================================================
\section{ICS-Inspired Intervention Implementation Plan}

The deployed intervention layer must not use surface normals, map geometry,
odometry, explicit velocity estimates, or privileged simulator state. It uses
MID360 ranges, masks, reliability weights, body-frame ray directions, command
velocity, and history.

\begin{longtable}{@{}L{0.27\textwidth}L{0.59\textwidth}@{}}
\caption{ICS-inspired attenuation details.}\\
\toprule
\textbf{Component} & \textbf{Plan} \\
\midrule
\endfirsthead
\toprule
\textbf{Component} & \textbf{Plan} \\
\midrule
\endhead
\(c_i(u)\) & For candidate command \(u=\vgov\), compute line-of-sight approach component along MID360 ray direction \(d_i\), weighted by \(m_iw_i\). \\
Range-rate fitting & Estimate \(\dot r_i\) from MID360 history and timestamps using robust finite differences or short linear fit. \\
\(\vsafe(d)\) & \(\vsafe(d)=\sqrt{2\amax\max(d-d_{\mathrm{safe}},0)}\). \\
Latency margin & Reduce effective clearance by command speed times sensor/controller latency plus configured margin. \\
Active set & Beams with valid, reliable returns and braking-feasibility risk. Invalid or stale beams are excluded or handled conservatively. \\
Empty active set & Return \(\beta_t=1\) and \(\vfin=\vgov\). \\
Emergency bypass & If reliable minimum clearance is below emergency threshold, ignore governor command and send \(\vbrake\)/hover command. \\
Logs & \(\beta_t\), active-set size, worst beam, minimum clearance, worst margin, emergency flag, \(\vgov\), \(\vfin\). \\
\bottomrule
\end{longtable}

Tests: monotonic \(\beta_t\) with speed/clearance, empty active set, emergency
bypass, no surface-normal imports/fields, no odometry/map access, and final
command bounds.

%============================================================
\section{Observability Logger Implementation Plan}

Range-Jacobian observability is an evaluation and analysis module, not a
deployed control dependency.

\paragraph{Simulation/offline logger.}
Use RayCaster scene access and simulator hit information to compute \(J\) from
surface normals if available, or from finite-difference ray casting. Surface
normals are allowed only in this offline logger. Log rank estimate,
\(\sigma_{\min}(J)\), condition number, weak direction, valid beam fraction,
scenario ID, and drift.

\paragraph{Hardware proxy.}
For real TASLAB\_UAV/MID360 trials, use scan geometry, valid-beam distribution,
range variation, and small commanded probes to compute an empirical
observability proxy. Label it as a proxy, not the exact \(J\).

Required plots: drift versus \(\sigma_{\min}(J)\), per-scenario drift ranking,
and weak-direction drift alignment.

%============================================================
\section{Reward and Training Plan}

Reward code currently lives inline in \code{NavigationEnv.\_compute\_state\_and\_obs}.
instinctRL should either refactor it to \code{instinctRL/rewards.py} or gate a
clearly marked instinctRL reward path in \code{env.py}. Privileged TASLAB state
may be used for reward and critic only.

\begin{longtable}{@{}L{0.28\textwidth}L{0.44\textwidth}L{0.18\textwidth}@{}}
\caption{Reward terms.}\\
\toprule
\textbf{Term} & \textbf{Plan} & \textbf{Privilege status} \\
\midrule
\endfirsthead
\toprule
\textbf{Term} & \textbf{Plan} & \textbf{Privilege status} \\
\midrule
\endhead
Tracking reward & Penalize unsafe deviation from \(\vcmd\) under nonzero command. & Actor-clean; actual velocity optional reward/eval only. \\
Anchor reward & Penalize masked MID360 anchor error under null command. & Actor-clean. \\
Safety penalty & Penalize low clearance and collision labels. & Reward/critic/eval only if from simulator. \\
Command-compliance gated by \(\beta_t\) & Do not punish attenuation when ICS marks command unsafe. & Actor-clean. \\
Intervention-usage penalty & Discourage policy reliance on ICS. & Uses ICS logs. \\
Smoothness penalty & Penalize changes in \(\vgov\) and \(\vfin\). & Actor-clean. \\
Collision penalty & Large terminal penalty from simulator/contact safety boundary. & Reward/eval only. \\
\bottomrule
\end{longtable}

%============================================================
\section{Baselines and Ablations}

\begin{longtable}{@{}L{0.08\textwidth}L{0.23\textwidth}L{0.25\textwidth}L{0.30\textwidth}@{}}
\caption{instinctRL baseline switches.}\\
\toprule
\textbf{ID} & \textbf{Config switch} & \textbf{Expected input fields} & \textbf{Contribution tested} \\
\midrule
\endfirsthead
\toprule
\textbf{ID} & \textbf{Config switch} & \textbf{Expected input fields} & \textbf{Contribution tested} \\
\midrule
\endhead
B0 & \code{instinctRL.baseline.id=direct\_velocity} & \(\vcmd\) only. & TASLAB velocity controller alone. \\
B1 & \code{state\_velocity\_upper\_bound} & Privileged state/velocity, labeled upper bound. & Gap to forbidden state-aided policy. \\
B2 & \code{no\_history} & Single-frame allowed MID360/IMU/command. & Value of history. \\
B3 & \code{no\_anchor} & Full actor inputs, no anchor loss/features. & Measurement-space anchor contribution. \\
B4 & \code{no\_imu} & MID360, command, previous output, history. & IMU cue contribution. \\
B5 & \code{no\_ics} & Full governor, \(\beta_t=1\). & ICS attenuation contribution. \\
B6 & \code{distance\_only\_filter} & Range minimum and command. & Distance-only filtering comparison. \\
B7 & \code{braking\_heuristic} & MID360 ranges, masks, command, braking params. & Nonlearned braking heuristic. \\
B8 & \code{simplified\_cbf\_vo\_optional} & Baseline-specific geometric inputs, not main method. & Optional filter comparison. \\
\bottomrule
\end{longtable}

Metrics for all relevant baselines: station-keeping drift, range-anchor error,
tracking RMSE, minimum clearance, collision rate, ICS violation rate, command
preservation ratio, intervention frequency, and observability metrics.

%============================================================
\section{Experiment Pipeline}

Required scenario groups are corner, flat wall, corridor/tunnel, open field,
cluttered indoor, unsafe command toward wall, gap/doorway, dynamic obstacle if
available, real-robot observable holding, and real-robot unsafe-command
approach. Scenarios may reuse \code{universal\_generator.py}, current obstacle
terrain, or dedicated test fixtures, but each run must log platform
\code{TASLAB\_UAV}, sensor \code{Livox MID360}, actor-input audit result, and
baseline ID.

The required metric set is station-keeping drift, range-anchor error, tracking
RMSE, minimum clearance, collision rate, ICS violation rate, command
preservation ratio, intervention frequency, observability metrics, and
real-robot repeatability. Pose/mocap/odom may be used for evaluation metrics
only and must be absent from actor inputs.

%============================================================
\section{Development Tickets}

\begin{longtable}{@{}L{0.14\textwidth}L{0.18\textwidth}L{0.38\textwidth}L{0.20\textwidth}@{}}
\caption{instinctRL implementation tickets.}\\
\toprule
\textbf{Ticket} & \textbf{Objective} & \textbf{Implementation tasks} & \textbf{Acceptance} \\
\midrule
\endfirsthead
\toprule
\textbf{Ticket} & \textbf{Objective} & \textbf{Implementation tasks} & \textbf{Acceptance} \\
\midrule
\endhead
instinctRL-0 & Platform and sensor infrastructure audit. & Verify TASLAB registry/config/USD/params/controller; wire MID360 to TASLAB base link; validate range vector, mask, weights, ray ordering, timestamps; run actor audit. & All audit criteria pass; no actor leak; no generic platform/sensor in normal configs. \\
instinctRL-A & Direct velocity-governor baseline. & Add instinctRL config namespace; create actor-clean command path; preserve TASLAB velocity-controller execution. & B0 runs on TASLAB/MID360; actor audit passes. \\
instinctRL-B & Observation/history buffer. & Build MID360 \(r,m,w\), timestamps, IMU cues, command, previous output, history. & Shape/timing/dropout tests pass. \\
instinctRL-C & Measurement-space anchor. & Add anchor manager, hysteresis, capture/reset, masked error, anchor reward/logs. & Anchor lifecycle tests pass; B3 ablation exists. \\
instinctRL-D & Observability logger. & Add offline range-Jacobian and hardware proxy logger, plots, drift correlation. & Required plots produced; no deployed normal use. \\
instinctRL-E & ICS-inspired attenuation. & Add range-rate fit, active set, \(\vsafe\), latency margin, \(\beta_t\), emergency bypass. & Monotonicity and no-normal tests pass; B5 exists. \\
instinctRL-F & Reward integration and training. & Add reward terms and critic-only privileged branch; prevent actor leakage. & Reward tests and first stable training run pass. \\
instinctRL-G & Baselines and ablations. & Add B0--B8 config isolation and evaluator matrix. & All baselines run/log required metrics. \\
instinctRL-H & Real-robot deployment validation. & Add MID360/IMU/command ROS wrapper, body-frame velocity output, bag replay audit. & No odom/map actor input; latency and safety logs present. \\
\bottomrule
\end{longtable}

\paragraph{Risk level and dependencies.}
instinctRL-0, B, C, E, F, and H are high risk. instinctRL-A, D, and G are
medium risk. All later tickets depend on instinctRL-0 because platform and
sensor lock violations invalidate results.

%============================================================
\section{Automated Tests and Audits}

\begin{itemize}
\item Platform lock: assert \code{cfg.drone.model\_name == "TaslabUAV"} and
sensor config identifies Livox MID360 in instinctRL runs.
\item MID360 range: stable ray count, ray ordering, frame convention, timestamp
monotonicity, mask/dropout handling, reliability bounds.
\item TASLAB velocity interface: corrected body-frame velocity command reaches
the existing \code{VelController}/\code{LeePositionController} adapter without
exposing motor/CTBR actions to the actor.
\item Actor input contract: no pose, odometry, explicit velocity, map, SLAM,
dynamic-obstacle privileged state, or root state in actor TensorDicts.
\item Observation history: fixed shapes, rollover, stale/dropped frame markers.
\item Anchor lifecycle: hysteresis, capture, reset, masked error, inactive loss.
\item ICS: monotonicity, empty active set \(\beta_t=1\), emergency bypass, no
surface-normal use in deployed code.
\item Rewards: each term activates under intended condition and privileged
quantities stay reward/critic/eval only.
\item Baselines: config isolation for B0--B8 and explicit input-schema logs.
\item Logging completeness: platform lock, actor audit, anchor, ICS,
observability, reward terms, and metrics emitted.
\end{itemize}

%============================================================
\section{Risks and Mitigation}

\begin{longtable}{@{}L{0.30\textwidth}L{0.58\textwidth}@{}}
\caption{instinctRL risks and mitigations.}\\
\toprule
\textbf{Risk} & \textbf{Mitigation} \\
\midrule
\endfirsthead
\toprule
\textbf{Risk} & \textbf{Mitigation} \\
\midrule
\endhead
Active RayCaster remains attached to Hummingbird path. & Make instinctRL-0 fail until the MID360 prim path resolves to TASLAB\_UAV base link. \\
MID360 helper exists but is not wired into training. & Reuse \code{create\_livox\_from\_hydra\_cfg} or integration helper; log active sensor source. \\
Low-level controller uses state internally. & State claim boundary: actor is pose/odometry-free; controller internals are conventional stabilization. \\
Open-field unobservability. & Report as method limitation using observability logger; do not hide degenerate cases. \\
Range dropout and latency. & Use masks, weights, timestamps, staleness, latency margin, and dropout tests. \\
Policy overuses ICS. & Penalize intervention usage and compare against B5 no-ICS. \\
Frame convention mismatch. & Test body-frame unit commands through TASLAB controller adapter and ROS wrapper. \\
Reward hacking. & Log component rewards and validate with baselines and held-out scenarios. \\
Sim-to-real gap. & Keep TASLAB parameters and MID360 noise/dropout; validate with bag replay and conservative hardware tests. \\
Surface normals leak into deployed safety. & Separate observability logger from deployed ICS and add static/runtime no-normal audits. \\
\bottomrule
\end{longtable}

%============================================================
\section{Final Deliverables}

The v1.1 platform-locked handbook is complete only if:
\begin{itemize}
\item all code/development names use \code{instinctRL};
\item MID360 and TASLAB\_UAV are locked as the normal platform/sensor;
\item instinctRL-0 audits MID360 range vector, mask, weights, ray ordering,
timestamps, and TASLAB velocity-command acceptance;
\item every instinctRL module maps to existing infrastructure first;
\item every proposed new file/function explains why existing infrastructure is
insufficient;
\item every deployed actor input is audited;
\item every forbidden shortcut has a test or manual check;
\item every scientific contribution has a supporting experiment;
\item any generic sensor/platform/simulator use is marked debugging-only.
\end{itemize}

\paragraph{Current unavoidable limitations to flag.}
The current repository is close to the locked platform, but not yet compliant:
\code{training/cfg/drone.yaml} selects \code{TaslabUAV}, while
\code{NavigationEnv} still hard-codes a Hummingbird LiDAR attachment path; the
active actor receives forbidden goal-relative pose and explicit velocity; and
the ROS wrapper uses odometry and map raycasting in the policy path. These are
implementation limitations, not method changes.

\end{document}
